\section{Problem Summary}
Reorder the numbers \(1, 2, \ldots, n\) so that every pair of neighboring values differs by something other than \(1\). If no such order exists, report that immediately.

\section{Editorial}
The small cases are the only impossible ones:
\begin{itemize}
\item \(n = 1\) is already valid
\item \(n = 2\) and \(n = 3\) have no solution
\end{itemize}

For \(n \ge 4\), a simple construction works: print all even numbers first, then all odd numbers.

Why does that help?
\begin{itemize}
\item inside the even block, consecutive numbers differ by \(2\)
\item inside the odd block, consecutive numbers also differ by \(2\)
\item the only boundary is from the largest even number to the smallest odd number, and for \(n \ge 4\) that gap is never \(1\)
\end{itemize}

So the entire permutation is valid.

\section{Complexity Analysis}
\begin{itemize}
\item Time: \(O(n)\)
\item Memory: \(O(1)\) besides the output itself
\end{itemize}

\section{Notes}
Do not overcomplicate the construction. The even-then-odd ordering is the intended simple pattern.
